\chapter{Các mô hình học sâu sử dụng cho NIDS}
\section{Mạng nơ-ron hồi quy – Recurrent Neural Network (RNN)}
RNN [9] được sử dụng do các đặc trưng mạnh mẽ của nó để tìm hiểu từ dữ liệu trước đó, sau đó điều chỉnh phản ứng của nó để dự đoán. Hơn nữa, nó kết hợp hai đặc trưng chính:
\begin{enumerate}
    \item Trạng thái ẩn phân tán cho phép chúng lưu trữ nhiều thông tin về quá khứ một cách hiệu quả. 
    \item Phi tuyến động cho phép chúng cập nhật trạng thái ẩn theo những cách phức tạp. 
\end{enumerate}

\subsection{Giới thiệu}
RNN, một phần mở rộng của mạng nơ-ron chuyển tiếp thẳng, được thiết kế để nhận dạng các mẫu trong chuỗi dữ liệu. Các RNN được gọi là hồi quy vì chúng thực hiện cùng một nhiệm vụ cho mọi mục của chuỗi với đầu ra phụ thuộc vào các tính toán trước đó [10]. 

\subsection{Kiến trúc RNN}
Bên trái của Hình 2.1 hiển thị ký hiệu của một RNN và ở bên phải một RNN đang được hủy cuộn (hoặc mở ra) thành một mạng đầy đủ. Bằng cách hủy cuộn, chúng tạo nên mạng cho chuỗi hoàn chỉnh.

\begin{table}
    \centering
    \caption{Caption}
    \begin{tabular}{cccc}
        1 & 2 & 3 & 4\\
        5 & 6 & 7 & 8\\
         &  &  & \\
         &  &  & \\
    \end{tabular}
    \label{tab:my_table}
\end{table}


\begin{table}[ht]
  \centering
  \caption{Hiệu năng trên UNSW-NB15 (tập test).}
  \label{tab:unswnb15-models}
  \begin{tabular}{l
                  S[table-format=1.3]
                  S[table-format=1.3]
                  S[table-format=1.3]
                  S[table-format=1.3]}
    \toprule
    {Mô hình} & {Accuracy} & {Precision} & {Recall} & {F1} \\
    \midrule
    Logistic Regression & 0.904 & 0.892 & 0.881 & 0.886 \\
    Random Forest       & 0.936 & 0.928 & 0.915 & 0.921 \\
    XGBoost             & 0.945 & 0.938 & 0.927 & 0.932 \\
    SVM (RBF)           & 0.925 & 0.912 & 0.906 & 0.909 \\
    CNN (1D)            & 0.952 & 0.943 & 0.938 & 0.940 \\
    \bottomrule
  \end{tabular}
\end{table}



% \enquote{Ví dụ dấu ngoặc kép}
